% F2 · El modelo de entidades de un conjunto de datos (D10·05)
\documentclass[border=10pt]{standalone}
% ------------------------------------------------------------------
% Preámbulo común de las figuras del curso C1
% Paleta y tipografía alineadas con ../../estilo.tex
% ------------------------------------------------------------------
\usepackage[utf8]{inputenc}
\usepackage[T1]{fontenc}
\usepackage[spanish,es-noquoting]{babel}
\usepackage{helvet}
\renewcommand{\familydefault}{\sfdefault}
\usepackage{xcolor}
\usepackage{tikz}
\usetikzlibrary{arrows.meta,positioning,decorations.pathreplacing,calc,shapes.geometric,fit,backgrounds}

\definecolor{cHead}{RGB}{68,84,106}
\definecolor{cAcc}{RGB}{0,102,153}
\definecolor{cGris}{gray}{0.42}
\definecolor{cSuave}{RGB}{240,244,248}
\definecolor{cBien}{RGB}{34,124,68}
\definecolor{cMal}{RGB}{176,42,42}
\definecolor{cBienSuave}{RGB}{230,243,234}
\definecolor{cMalSuave}{RGB}{250,235,233}
\definecolor{cAccSuave}{RGB}{225,238,246}

\tikzset{
  caja/.style     = {rectangle, rounded corners=2.5pt, draw=cHead, line width=.7pt,
                     fill=cSuave, align=center, inner sep=6pt, font=\small},
  cajaAcc/.style  = {caja, draw=cAcc, fill=cAccSuave},
  cajaOk/.style   = {caja, draw=cBien, fill=cBienSuave},
  cajaNo/.style   = {caja, draw=cMal, fill=cMalSuave, dashed},
  cajaTenue/.style= {caja, draw=cGris, fill=white},
  rotulo/.style   = {font=\scriptsize\itshape, text=cGris, align=center},
  etiq/.style     = {font=\scriptsize, text=cHead, align=center,
                     fill=white, inner sep=1.5pt},
  fl/.style       = {-{Stealth[length=4.5pt]}, draw=cHead, line width=.7pt},
  flAcc/.style    = {fl, draw=cAcc},
  flNo/.style     = {-{Stealth[length=4.5pt]}, draw=cMal, line width=.7pt, dashed},
  llave/.style    = {decorate, decoration={brace, amplitude=4pt, raise=2pt},
                     draw=cGris, line width=.6pt},
  titulo/.style   = {font=\small\bfseries, text=cHead},
}

% Cabecera de figura: #1 número, #2 título
\newcommand{\figtit}[2]{\node[font=\footnotesize\bfseries, text=cGris, anchor=west]}

\begin{document}
\begin{tikzpicture}

\tikzset{
  ent/.style = {caja, text width=3.7cm, align=left, inner sep=7pt, font=\scriptsize},
  entAncha/.style = {ent, text width=4.3cm},
  card/.style= {etiq, font=\scriptsize\bfseries, text=cAcc},
}

% #1 nombre  #2 «una fila es»  #3 clave  #4 filas
\newcommand{\entidad}[4]{%
  {\small\bfseries\ttfamily #1}\\[2pt]
  \textcolor{cGris}{una fila es}\\ #2\\[3pt]
  \textcolor{cGris}{clave:} \texttt{#3}\\
  \textcolor{cGris}{filas:} #4}

\node[ent] (pac) at (0,0)
  {\entidad{pacientes}{un paciente incluido}{id\_paciente}{400}};

\node[ent] (vis) at (6.4,2.35)
  {\entidad{visitas}{una visita de un paciente}{id\_paciente + n\_visita}{$\sim$1.800}};

\node[ent] (mue) at (6.4,-0.35)
  {\entidad{muestras}{una muestra extraída}{id\_muestra}{$\sim$900}};

\node[ent] (img) at (6.4,-2.95)
  {\entidad{imagen}{un estudio de resonancia}{id\_estudio}{120}};

\node[entAncha] (det) at (12.6,-0.35)
  {\entidad{determinaciones}{un resultado de laboratorio}{id\_muestra + cod\_prueba}{$\sim$5.400}};

\draw[flAcc] (pac.east) -- (vis.west) node[card, pos=.55] {1 $\rightarrow$ 0 a 5};
\draw[flAcc] (pac.east) -- (mue.west) node[card, pos=.55] {1 $\rightarrow$ 0 a 3};
\draw[flAcc] (pac.east) -- (img.west) node[card, pos=.55] {1 $\rightarrow$ 0 o 1};
\draw[flAcc] (mue.east) -- (det.west) node[card, pos=.5]  {1 $\rightarrow$ n};

\node[rotulo, anchor=north west, text width=14.5cm, align=left] at (0,-4.55)
  {Ni una palabra de bases de datos: qué tablas hay, qué representa una fila en cada una,
   cuántas hay por paciente y por qué campo se enlazan. El diccionario de variables
   cuelga de cada tabla, no de un saco único.};

\end{tikzpicture}
\end{document}
